\scanpage{80}
% archive sheet number: 79

\deleted{deux adjoints de \(f\) existent et \unsure{soient|sont} isom. Il suffirait de l'avoir pour \deleted{des} catégories \deleted{\unsure{non vides}} \(C,C'\), telles que \(f,g\) ne soient pas des \deleted{\unsure{équivalences|foncteurs}}\inserted{pl. fid.} (on passe ensuite à \(\widehat C,\widehat C'\)).}

\deleted{Peut-on \(=\) prendre \(C=C', f=g\), i.e. peut on trouver une catégorie \(C\), et un endom \(f:C\unsure{-|\longrightarrow}C\) qui ne soit pas \deleted{une équivalence}\inserted{pl. fid.}, et qui soit isom. à son adjoint (à deux \unsure{bords|côtés})?}

OK. On soit \(C\) une catégorie ayant un objet nul \deleted{\unsure{0}} \(e\), mais tel que le foncteur constant \deleted{\unsure{\(C\longrightarrow\)}} \(C\longrightarrow C\) défini par \(f(x)=e\) ne soit pas pl fidèle i.e. \(C\) non une catégorie \unsure{discrétable|discrète|discrètisable}, donc \(C\) contient un objet \(x\unsure{\ne|=}e\), [Ex. minimal :

\diagram{\begin{tikzcd}[ampersand replacement=\&]e\arrow[r,shift right,"v"']\&a\arrow[l,shift right,"u"']\arrow[loop right,"{\deleted{\unsure{vu}}w}"]\end{tikzcd}\,:}

deux objets \(e,a\), trois flèches en plus des flèches identiques \(u,v,vu=w\), relations \deleted{\unsure{à prio.}}\inserted{de} composition sont déterminées \unsure{uniqt|uniquement} alors par ces conditions, \deleted{sauf} une [\(uv=\operatorname{id}_e\quad uw=u,wv=v\)], savoir que il faut dire ce qu'est \(w^2\) : On peut prendre \(w^2=\operatorname{id}_a\) ou \(w^2=w\) au choix], On trouve le
